\chapter{Derivational morphology}
\label{chap:derivations}

\section{Stress in derivation}
\label{sec:affix-stress}

Derivational suffixes belong to three stress classes. A stress-neutral suffix (Class W) preserves the stress of the base. A secondary-stressed suffix (Class S) preserves the base's primary stress and adds secondary stress of its own. A stress-attracting suffix (Class A) assigns a new primary-stress position. These classes are shown in Table \ref{tab:affix-stress}. Ordinary inflectional endings are stress-neutral.

\begin{table}[!ht]
    \begin{tabularx}{\textwidth}{l X X}
        Class & Stress behaviour & Suffixes \\
        W & Retains the stress of the base & \nk{-ing/-ning}, \nk{-æg/-gæg}, \nk{-inn/-ginn}, \nk{-ats} \\
        S & Retains primary stress and adds secondary stress & \nk{-hæd}, \nk{-bägd} \\
        A & Attracts primary stress & \nk{-tet/-itet}, \nk{-iser/-tiser} \\
    \end{tabularx}
    \caption{Stress classes of derivational suffixes}
    \label{tab:affix-stress}
\end{table}

In the stress-only notation described in Chapter \ref{chap:phonology}, \S\ref{sec:stress-prosody}, \nk{ˈhond} becomes \nk{ˈhonding}, retaining its initial stress, while \nk{ˈgrom} becomes \nk{ˈgromˌhæd}, with secondary stress on the suffix. In \nk{bruˈtal} $\rightarrow$ \nk{brutaliˈtet}, primary stress moves to the final syllable. The stress marks indicate prominence without specifying vowel or consonant quantity and are not part of ordinary spelling.

\section{Survey of derivational affixes}

\subsection{-ats}

The suffix \nk{-ats} is used to form nouns from verbs ending in \nk{-iser} (see \S\ref{sec:suffix-iser}). The resulting noun denotes the action of the verb and is feminine. The suffix is stress-neutral: \nk{normaˈliser} $\rightarrow$ \nk{normaˈliserats}.

\begin{itemize}
    \item \nikp{dramatiser}{to dramatize} $\rightarrow$ \nikp{dramatiserats}{dramatization}
    \item \nikp{kannibaliser}{to cannibalize} $\rightarrow$ \nikp{kannibaliserats}{cannibalization}
    \item \nikp{normaliser}{to normalize} $\rightarrow$ \nikp{normaliserats}{normalization}
\end{itemize}

\subsection{-bägd}

The suffix \nk{-bägd} is used to form adjectives from nouns and verbs. The resulting adjective denotes the ability to perform the action denoted by the verb or the quality of the noun.

The base retains its primary stress, while \nk{-bägd} receives secondary stress, as in \nk{ˈvædˌbägd}. The secondary prominence is clearest in careful speech.

\begin{itemize}
    \item \nikp{vædbägd}{visible} $\leftarrow$ \nikp{væd}{see}
    \item \nikp{kanselerbägd}{cancellable} $\leftarrow$ \nikp{kanseler}{cancel}
    \item \nikp{klikkerbägd}{clickable} $\leftarrow$ \nikp{klikker}{click}
    \item 
\end{itemize}

\subsection{be-}


\subsection{-hæd}

The suffix \nk{-hæd} is used to form abstract nouns from adjectives and some nouns. The resulting nouns takes the meaning of the adjective or noun and turns it into an abstract concept. The suffix is productive with adjectives but less so with nouns. The resulting noun takes the feminine gender.

The suffix receives secondary stress without moving the base's primary stress, as in \nk{ˈgromˌhæd}. As with \nk{-bägd}, this secondary prominence is clearest in careful speech.

\begin{itemize}
    \item \nikp{grom}{stupid} $\rightarrow$ \nikp{gromhæd}{stupidity}
    \item \nikp{jærv}{tired} $\rightarrow$ \nikp{jærvhæd}{tiredness}
    \item \nikp{maus}{happy} $\rightarrow$ \nikp{maushæd}{happiness}
    \item \nikp{ngord}{sibling} $\rightarrow$ \nikp{ngordhæd}{brotherhood/sisterhood}
\end{itemize}

\subsection{-ing}
\label{sec:diminutives}

The suffix \nk{-ing} is used to form diminutives from nouns. The resulting noun takes the neuter gender. After nouns ending in a vowel or in \nk{s}, as well as some exceptions, the suffix is \nk{-ning}.

Both forms of the suffix are stress-neutral and retain the lexical stress of the noun, as in \nk{ˈhond} $\rightarrow$ \nk{ˈhonding}.

\begin{itemize}
    \item \nikp{stor}{brother} $\rightarrow$ \nikp{storning}{brother, little brother}
    \item \nikp{hond}{child} $\rightarrow$ \nikp{honding}{child, little child}
\end{itemize}

\subsection{-iser}

\label{sec:suffix-iser}

The suffix \nk{-iser} is used to form verbs from nouns. The resulting verb denotes the action of the noun. The suffix is productive with nouns of foreign (mainly Greek, entering Niklorsk through Latin, French or English) origin. After Greek roots ending in \nk{-a} the suffix is \nk{-tiser}. Some loanwords ending in \nk{-iser/-tiser} do not have a corresponding nativized noun.

Primary stress falls on the syllable containing the first vowel of \nk{-iser/-tiser}, as in \nk{normaˈliser} and \nk{dramaˈtiser}. This stressed vowel remains when the final \nk{-er} is replaced or deleted in conjugation: \nk{normaˈliser} $\rightarrow$ \nk{normaˈlisatt}. The inflectional ending does not move the stress (see Chapter \ref{chap:verbs}, Table \ref{tab:verb-conjugation-latinate}).

\begin{itemize}
    \item \nikp{drama}{drama} $\rightarrow$ \nikp{dramatiser}{to dramatize}
    \item \nikp{kannibal}{cannibal} $\rightarrow$ \nikp{kannibaliser}{to cannibalize}
    \item \nikp{normal}{normal} $\rightarrow$ \nikp{normaliser}{to normalize}
\end{itemize}

With defective roots:

\begin{itemize}
    \item \nikp{organiser}{organize}
\end{itemize}

\subsection{-itet}

See \S\ref{sec:suffix-tet}.

\subsection{-ning}

See \S\ref{sec:diminutives}.

\subsection{-tet}
\label{sec:suffix-tet}

The suffix \nk{-tet/-itet} is used to form noun from adjectives of foreign origin. The resulting noun denotes the quality of the adjective and is feminine.

Both forms of the suffix attract primary stress to \nk{-tet}, as in \nk{brutaliˈtet} and \nk{objektiviˈtet}. Subsequent inflection preserves this stress.

\begin{itemize}
    \item \nikp{passiv}{passive} $\rightarrow$ \nikp{passivetet}{passivity}
    \item \nikp{brutal}{brutal} $\rightarrow$ \nikp{brutalitet}{brutality}
    \item \nikp{objektiv}{objective} $\rightarrow$ \nikp{objektivitet}{objectivity}
\end{itemize}


\section{Compound words}

\subsection{Noun-noun compounds}
\label{sec:compound-stress}
Two nouns can be combined to form a compound noun. In the orthography, the resulting word is written as a single word. The second noun is the head of the compound and the first noun serves as a modifier. The compound noun takes the gender of the head noun. Morphologically, noun-noun compounds are formed by juxtaposing the two nouns without any intervening morphemes, although a small number of compounds have an epenthetic \nk{-e-} or \nk{-en-}, in case the resulting compound word results in an illegal consonant cluster.

The first constituent normally carries primary stress, while the second carries secondary stress. Each prominence falls on the constituent's lexical stressed syllable. Thus a first constituent with non-initial lexical stress retains that position rather than shifting stress to the beginning of the compound. The epenthetic linking elements \nk{-e-} and \nk{-en-} remain unstressed. Secondary stress preserves vowel clarity and some duration in the second constituent (see Chapter \ref{chap:phonology}, \S\ref{sec:stress-prosody}).

The following compounds have the stress patterns \nk{ˈplastˌrum}, \nk{ˈsnörˌhus}, and \nk{ˈrörˌdygg} in stress-only notation.

\noindent\begin{tblr}{colspec = {XX}}
    \nikp{plast}{plastic} + \nikp{rum}{utensil} & \nikp{plastrum}{plastic utensil} \\
    \nikp{snör}{beach} + \nikp{hus}{house} & \nikp{snörhus}{beach house} \\
    \nikp{rör}{rye} + \nikp{dygg}{bread} & \nikp{rördygg}{rye bread} \\
\end{tblr}
